\section{Week 3 Refresher: Tokenization and Embedding Geometry}

\subsection*{Setting the Scene}
This week answers a hidden question: what mathematical object is a ``token'' once the model starts computing? The field evolved from count-based lexical models to distributed vectors, then to context-dependent embeddings in transformer systems. The critical insight was that representation is not a cosmetic preprocessing step. It defines the geometry in which similarity, retrieval, and attention operate.

\subsection*{Notation Introduced in This Chapter}
\begin{itemize}[leftmargin=1.5em]
  \item $V$: tokenizer vocabulary (finite symbol set).
  \item $\iota(\cdot)$: tokenizer mapping from text to token IDs.
  \item $\mathrm{onehot}(i)$: basis vector for token index $i$.
  \item $E$: embedding matrix.
  \item $e_i$: embedding vector for token $i$.
  \item $u,v,q,d$: generic vectors used in similarity comparisons.
\end{itemize}

\subsection*{What You Need To Recall Precisely}
\begin{itemize}[leftmargin=1.5em]
  \item \textbf{Tokenizer as function.} A tokenizer defines mapping $\iota:\text{text}\to\{1,\dots,|V|\}$. Every downstream probability is conditioned on this representation choice.
  \item \textbf{One-hot basis and embedding lookup.}
  \[
    e_i = E^\top \mathrm{onehot}(i).
  \]
  This is linear algebra row selection.
  \item \textbf{Segmentation criteria.} BPE, WordPiece, and Unigram optimize different objectives (merge frequency, likelihood-style score, probabilistic segmentation).
  \item \textbf{Similarity metrics.}
  \[
    u^\top v \quad \text{vs} \quad \cos(u,v)=\frac{u^\top v}{\|u\|\|v\|}.
  \]
  Dot product includes magnitude effects; cosine removes scale.
  \item \textbf{Anisotropy and concentration.} Contextual embeddings can cluster in narrow directional cones, weakening naive nearest-neighbor interpretations.
\end{itemize}

\subsection*{How These Results Build on Each Other}
\begin{enumerate}[leftmargin=1.5em]
  \item Tokenization defines discrete symbols.
  \item Symbol IDs become one-hot basis vectors.
  \item Embedding matrix maps one-hot vectors into continuous space.
  \item Similarity metric defines neighborhood structure.
  \item Neighborhood structure influences retrieval and downstream context selection.
\end{enumerate}

\subsection*{Reminder Glossary (Term by Term)}
\begin{itemize}[leftmargin=1.5em]
  \item \textbf{Vocabulary:} finite symbol inventory used by the model.
  \item \textbf{Token ID:} integer index assigned by tokenizer.
  \item \textbf{One-hot vector:} canonical basis vector for a token ID.
  \item \textbf{Embedding matrix:} learned lookup table from IDs to vectors.
  \item \textbf{Dot product:} magnitude-sensitive alignment score.
  \item \textbf{Cosine similarity:} direction-only alignment score.
  \item \textbf{Anisotropy:} directional imbalance in vector space.
\end{itemize}

\subsection*{Worked Example (Why Metric Choice Changes Outcomes)}
Let query vector be $q=(10,0)$ and two candidates:
\[
d_1=(1,0),\quad d_2=(1,1).
\]
Dot product ranking:
\[
q^\top d_1 = 10,\quad q^\top d_2 = 10.
\]
Tie under dot product.
Cosine ranking:
\[
\cos(q,d_1)=1,\quad \cos(q,d_2)=\frac{1}{\sqrt{2}}\approx0.707.
\]
Now $d_1$ is clearly preferred directionally.  
Reasoning implication: if your retrieval objective is semantic direction, cosine normalization is often essential; if magnitude encodes useful confidence/strength, dot-product behavior may be intentional.

\subsection*{Intuition Checkpoints}
\begin{itemize}[leftmargin=1.5em]
  \item Keep sequence-position indices separate from vocabulary indices.
  \item Tokenization choice changes sequence length and therefore perplexity comparability.
  \item Static-word-vector intuitions transfer only partially to contextual embeddings.
\end{itemize}

\subsection*{Pointers}
\begin{itemize}[leftmargin=1.5em]
  \item Probability/perplexity interpretation: see Week 1 refresher.
\end{itemize}

\subsection*{References and What To Use Them For}
\begin{itemize}[leftmargin=1.5em]
  \item \textbf{Jurafsky--Martin, \emph{SLP}} --- tokenizer design and language representation.
  \item \textbf{Strang, \emph{Linear Algebra}} --- vector spaces, norms, geometry.
  \item \textbf{Mikolov et al.} --- distributed embedding intuition.
  \item \textbf{CS224N materials} --- modern embedding/retrieval implementation context.
\end{itemize}
